\begin{lemma}[Smooth1Properties-below]
\uses{Smooth1}
\label{Smooth1Properties_below}
\lean{Smooth1Properties_below}
% at /home/runner/work/trough/trough/library/PrimeNumberTheoremAnd/MellinCalculus.lean:886.0-957.61
\leanok
Fix $\epsilon>0$. There is an absolute constant $c>0$ so that:
  If $0 < x \leq (1-c\epsilon)$, then
  $$\widetilde{1_{\epsilon}}(x) = 1.$$
\end{lemma}
\begin{proof}
\uses{DeltaSpikeMass,DeltaSpike,MellinConvolution}
\leanok
Opening the definition, we have that the Mellin convolution of $1_{(0,1]}$ with $\nu_\epsilon$ is
  $$
  \int_0^\infty 1_{(0,1]}(y)\nu_\epsilon(x/y)\frac{dy}{y}
  =
  \int_0^1 \nu_\epsilon(x/y)\frac{dy}{y}.
  $$
  The support of $\nu_\epsilon$ is contained in $[1/2^\epsilon,2^\epsilon]$, so it suffices to
  consider $y \in [1/2^\epsilon x,2^\epsilon x]$ for nonzero contributions. If $x < 2^{-\epsilon}$,
  then the integral is the same as that over $(0,\infty)$:
  $$
  \int_0^1 \nu_\epsilon(x/y)\frac{dy}{y}
  =
  \int_0^\infty \nu_\epsilon(x/y)\frac{dy}{y},
  $$
  in which we change variables to $z=x/y$ (using $x>0$):
  $$
  \int_0^\infty \nu_\epsilon(x/y)\frac{dy}{y}
  =
  \int_0^\infty \nu_\epsilon(z)\frac{dz}{z},
  $$
  which is equal to one by Lemma \ref{DeltaSpikeMass}.
  We then choose
  $$
    c:=\log 2,
  $$
  which satisfies
  $$
    c > \frac{1-2^{-\epsilon}}\epsilon
  $$
  by Lemma \ref{Smooth1Properties_estimate}, so
  $$
    1-c\epsilon < 2^{-\epsilon}.
  $$
\end{proof}
